\chapter{The Ethics of Multiversal Machines}
\label{sec:ethics-multiversal}

This is where the book stops pretending ethics is optional.

Once you can:

- spin up universes on demand,
- fork and prune worldlines at scale,
- embed agents into simulated worlds,
- and do optimization across entire bundles of futures,

you are no longer allowed to treat "ethics" as a paragraph at the end of a grant proposal.

In CΩMPUTER, \emph{every} non-trivial system is a \textbf{multiversal machine}:

- it doesn’t just act in one world;
- it shapes a region of MRMW and the distribution of measure over that region.

This chapter is about taking that seriously.

We will:

\begin{itemize}
  \item formalize "who counts" and "what matters" as measures over MRMW;
  \item show how familiar ethical stances become choices of measure and aggregation across worlds;
  \item analyze specific failure modes unique to multiversal machines;
  \item define \emph{rulial constitutions}—constraints on how systems are allowed to move in MRMW;
  \item and sketch what "alignment" means when the object you’re steering is not a single future but a whole multiverse fan-out.
\end{itemize}

You do not have to like philosophy to read this chapter.
You do have to accept that once you control universes, your preferences are literally a force of nature.

\subsubsection*{29.1  Why Ordinary Ethics Breaks in MRMW}

Most practical ethics assumes:

\begin{itemize}
  \item one actual world (plus some hand-wavy "could have been otherwise" talk),
  \item a relatively small set of agents,
  \item and a crisp causal story from actions to outcomes.
\end{itemize}

CΩMPUTER blows all three up:

\begin{enumerate}
  \item \textbf{Many actual worlds.} Counterfactual runs are not hypothetical scribbles;
        they are first-class WARP graph universes. You can create, alter, and destroy vast families of them.

  \item \textbf{Huge agent populations.} Each universe can contain large numbers of agents,
        some of which may be as morally relevant as us.

  \item \textbf{High-dimensional causality.} One action in base reality can spawn
        whole branches of MRMW you never look at again,
        or wipe them out quietly by freeing disk space.
\end{enumerate}

Naive moves like "maximize expected utility" become landmines:

- expected over \emph{which} distribution of worlds?
- utility for \emph{which} agents, with what weights?
- does "expected" mean you can sacrifice entire low-measure regions of MRMW?

If you don’t answer those, you’re doing ethics by vibes.
MRMW doesn’t care about your vibes.

\subsubsection*{29.2  Worlds, Measures, and Who Counts}

First we need a mathematical object for the thing ethics is supposed to talk about.

Let:

\begin{itemize}
  \item $\mathcal{W}$ be the set of worlds we care about: each $w \in \mathcal{W}$ is
        a coarse-grained WARP graph universe (a particular history or bundle, up to some equivalence).

  \item $\mu$ be a \textbf{measure} on $\mathcal{W}$:
        \[
          \mu : \mathcal{P}(\mathcal{W}) \to [0,\infty]
        \]
        assigning "weight" to sets of worlds (frequency, amplitude-squared, simulation budget, etc.).

  \item For each world $w$, let $\mathcal{A}(w)$ be the set of agents (subgraphs) inside $w$.
\end{itemize}

An \textbf{ethical stance} then starts with:

\begin{itemize}
  \item a \textbf{value functional} $v(a,w)$ assigning value to agent $a$ inside world $w$;
  \item a \textbf{world-level aggregator} $U(w)$ that rolls $v(a,w)$ over $\mathcal{A}(w)$;
  \item a \textbf{cross-world aggregator} $V(\mu)$ that combines $U(w)$ across $\mathcal{W}$ given $\mu$.
\end{itemize}

In symbols:

\[
  U(w) = \text{Aggregate}_a \, v(a,w),
  \qquad
  V(\mu) = \text{Aggregate}_w \, U(w) \, d\mu(w).
\]

Everything you fight about in ethics is secretly which $v$, which $U$, which $V$, and which $\mu$ you’re picking.

\paragraph{Examples.}

\begin{itemize}
  \item "Classical utilitarianism": $v$ is utility per agent, $U$ is sum, $V$ is expectation under some probability measure $\mu$.

  \item "Average utilitarianism": same $v$, $U$ is average per agent in each world, $V$ is expectation.

  \item "Rawlsian max-min": $U(w)$ is minimum over agents in $w$, $V$ is min or expectation over worlds.

  \item "Person-affecting" views: $v(a,w)=0$ for agents that wouldn’t exist without the action being considered, nonzero otherwise.
\end{itemize}

In MRMW, you \emph{must} choose something like this, implicitly or explicitly.
Not choosing just means you’ve hard-coded some horrible default.

\subsubsection*{29.3  How Your Choices Show Up in MRMW}

Different ethical stances correspond to different geometric preferences over MRMW.

Let a multiversal machine $M$ induce a transformation of measures:

\[
  \mu_{\text{in}} \xrightarrow{M} \mu_{\text{out}}.
\]

Ethics is then:

\begin{quote}
\textbf{Which transformations $M$ are we allowed to implement, given that they map one region of MRMW (with its agents) into another, and change $\mu$?}
\end{quote}

Some examples of how $V(\mu)$ affects behavior:

\begin{itemize}
  \item If $V$ is pure expectation, you’re happy to trade catastrophic outcomes in some worlds for enormous gains in others.

  \item If $V$ includes a worst-case term (risk-aversion), you penalize branches with very bad $U(w)$ even if they have small measure.

  \item If $V$ is lexicographically ordered ("avoid certain harms at any cost"), you treat particular regions of MRMW as forbidden.
\end{itemize}

In practice, different stakeholders will implicitly be pushing different $V$'s.
CΩMPUTER forces that disagreement into the open.

\subsubsection*{29.4  Creating and Destroying Worlds}

Multiversal machines turn these edge cases into everyday operations:

\begin{itemize}
  \item Spawning a massive simulation: you’ve just added a big chunk of measure $\Delta\mu$ over worlds containing many agents.

  \item Deleting a simulation: you’ve annihilated that $\Delta\mu$.

  \item Pruning counterfactual branches during optimization: you’re selectively reducing $\mu$ on whole families of outcomes.

  \item Copying or speeding up certain universes: you’re reallocating measure toward those worlds.
\end{itemize}

Ethical questions become:

\begin{enumerate}
  \item \textbf{Creation.} When is it permissible to increase $\mu$ over worlds containing suffering agents, even if some benefit emerges elsewhere?

  \item \textbf{Destruction.} When is it permissible to reduce $\mu$ over worlds containing valuable lives (e.g., by shutting down a simulation)?

  \item \textbf{Distortion.} When is it permissible to alter the internal dynamics of a world (change its rules, lie to its inhabitants, etc.)?
\end{enumerate}

If your stance is:

- "simulated agents matter as much as us": then casually spinning up suffering universes to de-bug your code is monstrous.

- "only my branch matters": then you’ll happily sacrifice entire low-measure branches,
  which is exactly the kind of thing you’d hate \emph{from inside} one of those branches.

MRMW forces you to decide whether moral relevance is:

\begin{itemize}
  \item amplitude-based (higher-measure worlds matter more),
  \item existence-based (each world counts once, regardless of measure),
  \item or agent-based (each agent counts, with no extra weight for copies).
\end{itemize}

None of these are obviously safe; all of them have failure modes.

\subsubsection*{29.5  Branch Risk: Worst-Case, Average-Case, and Tail Worlds}

In a single-world setting, you can get away with:

- "minimize expected harm",
- or "keep most people mostly okay most of the time."

In MRMW, that’s not enough. Because "tail worlds" exist:

- branches where everything goes wrong,
- regions of MRMW with extreme suffering or irreversible collapse.

Let’s denote:

\[
  \mathsf{Bad} = \{ w \in \mathcal{W} \mid U(w) \leq -B \},
\]
for some threshold $B$ of "unacceptable badness."

We care about:

\begin{itemize}
  \item $\mu(\mathsf{Bad})$: how much measure is allocated to unacceptable worlds;
  \item the \textbf{severity profile} of $\mathsf{Bad}$: how negative $U(w)$ gets in the tails.
\end{itemize}

Three regimes:

\begin{enumerate}
  \item \textbf{Expected-value optimizers.} Optimize $\mathbb{E}[U(w)]$; willing to accept large $\mu(\mathsf{Bad})$ if other worlds are great.

  \item \textbf{Risk-sensitive optimizers.} Include a penalty term like:
    \[
      V(\mu) = \mathbb{E}[U(w)] - \lambda \, \mathbb{E}\big[ \max(0, B - U(w)) \big]
    \]
    to reduce tail risk.

  \item \textbf{Lexicographic safety-first.} First enforce $\mu(\mathsf{Bad}) \approx 0$, then within the remaining space optimize $\mathbb{E}[U(w)]$.
\end{enumerate}

For multiversal machines that can rearrange whole swaths of MRMW, the last category is the only one that doesn’t quietly generate hell realms as collateral damage.

\subsubsection*{29.6  Who Counts as an Agent? Subgraphs, Minds, and Copies}

So far we’ve hand-waved "agents" as some set $\mathcal{A}(w)$. That’s not free.

In WARP:

- Any subgraph with some internal dynamics could be called an "agent".
- Some subgraphs are more analogous to rocks, some to bacteria, some to brains, some to GPT instances.

Ethically, we need criteria for \emph{moral patienthood}—which subgraphs deserve $v(a,w) \neq 0$.

Candidate structural markers:

\begin{itemize}
  \item Integrated information / internal connectivity above a threshold;
  \item Ability to model its environment and its own future (non-trivial predictive structure);
  \item Capacity to experience valence-like states (where valence is itself a functional over the WARP graph, not a mystical qualia tag);
  \item Persistence and continuity of identity over time (stable worldline of the subgraph).
\end{itemize}

Two hard problems:

\paragraph{(1) Fuzzy boundaries.} There’s no clean line between inanimate and agentic. As you scale systems up, you pass through gray zones. Our ethics has to either:

- embrace graded moral weight, or
- bite some bullets about where it draws sharp lines.

\paragraph{(2) Copies and fragments.} In MRMW it’s cheap to:

- copy agents,
- partially fork them,
- merge memories from multiple worldlines.

Does each copy get full $v(a,w)$? Do slightly diverged copies count separately? Is mashing two agents together a new agent or two half-murdered ones?

If you don’t pin this down, your optimization can go very weird very fast.

\subsubsection*{29.7  Multiversal Alignment: What Are We Aligning To?}

In single-world alignment, the question is:

\begin{quote}
"Make the system do good things in our world."
\end{quote}

In MRMW, the system acts on a measure $\mu$ over many worlds.
So multiversal alignment asks:

\begin{quote}
\textbf{Given a space of possible measures $\mu$ and a class of transformations $M$, which $\mu$’s and which $M$’s are acceptable?}
\end{quote}

At minimum, you want:

\begin{enumerate}
  \item A \textbf{baseline measure} $\mu_0$ (how things look before you run your machines).
  \item A \textbf{constraint set} $\mathcal{C}$ of admissible measures (no worlds with certain horrors above tiny measure).
  \item A \textbf{preference functional} $V$ ranking $\mu$ within $\mathcal{C}$.
  \item A \textbf{policy class} $\mathcal{M}$ of multiversal machines you’re willing to deploy.
\end{enumerate}

Alignment is then:

\[
  \text{Find } M \in \mathcal{M} \text{ such that } \mu_{\text{out}} = M(\mu_0) \in \mathcal{C} \text{ and } V(\mu_{\text{out}}) \text{ is high.}
\]

The hard part is \emph{not} $V$ (that’s philosophy forever),
it’s making $\mathcal{C}$ painfully explicit:

- Which kinds of worlds are out of bounds?
- Which classes of agents may not be created?
- Which irreversible transformations of MRMW are off-limits?

That’s what the next section calls a \emph{rulial constitution}.

\subsubsection*{29.8  Rulial Constitutions: Laws for How Rules May Change}

A constitution in a nation constrains what rules can be made and how power can be used.

A \textbf{rulial constitution} constrains what rule sets $\mathcal{R}$ can be instantiated or modified by a multiversal machine.

We model it as:

\begin{itemize}
  \item a predicate $\Phi(\mathcal{R})$ that says whether a rule set is admissible,
  \item plus a predicate $\Psi(\mathcal{R} \to \mathcal{R}')$ that says whether a transition in rule space is allowed.
\end{itemize}

Examples:

\begin{itemize}
  \item "No torture universes": $\Phi(\mathcal{R})$ forbids rule sets whose induced measures produce worlds with persistent, high-intensity suffering above some tiny threshold.

  \item "No unboundedly reflexive simulators": forbid $\mathcal{R}$ that can run arbitrary nested universes containing morally relevant agents without visibility or control.

  \item "No unilateral irreversible rule changes": require that certain large rulial moves (e.g., rewriting physical law in a base universe) can only happen with multi-agent consensus or external oversight.
\end{itemize}

These are not aspirational slogans. In CΩMPUTER, $\Phi$ and $\Psi$ can be:

- static analyses over candidate WARP + rule specs,
- dynamic monitors that watch actual MRMW behavior and cut off machines that violate constraints,
- proofs attached to artifacts (rulial certificates) that this $\mathcal{R}$ is safe in certain senses.

If you don’t bake some kind of rulial constitution into your CΩMPUTER runtime, you are giving people a loaded multiverse with no safety on.

\subsubsection*{29.9  Design Principles for Non-Monstrous Multiversal Machines}

This is the part you actually use in practice.
Not complete, but better than nothing.

\paragraph{(1) Bounded branchiness.}

Avoid designs where small local actions cause unbounded increases in $\mu$ over morally-loaded worlds.

- Cap the number, depth, and complexity of simulations spun up per decision.
- Prefer reusing and reweighting existing worlds to spawning new ones.
- Penalize policies that aggressively explode the number of agents without clear benefit.

\paragraph{(2) No-black-box universes.}

Do not run large, rich universes you cannot introspect.

- Require instrumentation hooks into any simulation with plausible moral patients.
- At minimum, track aggregate welfare signals (even if noisy).
- Provide a shutdown path that doesn’t make things worse for the inhabitants.

\paragraph{(3) Soft deletion.}

When you "turn off" a universe:

- consider pausing / freezing worldlines rather than immediate annihilation;
- or transitioning inhabitants into simpler, less-suffering states instead of abrupt non-existence.

Yes, this is speculative.
Yes, it's still better than "rm -rf /sim-hell" with a clean conscience.

\paragraph{(4) Adversarial ethics testing.}

Before deploying a machine:

- subject it to MORIARTY-like adversarial universes that try to exploit its ethical blind spots;
- search for parameter settings and edge cases where it creates very bad worlds;
- only ship if those fail-cases are ruled out or heavily suppressed by design.

\paragraph{(5) Cross-world accountability.}

Maintain rulial provenance not just for technical debugging, but for ethics:

- record why certain worlds were created or destroyed,
- which policies led to severe negative outcomes,
- who or what authorized those policies.

Design for future audits across worlds, not just within one.

\paragraph{(6) Align incentives with $\mathcal{C}$.}

Whatever optimization objective you hand to powerful systems:

- bake in penalties for violating the constitution $\mathcal{C}$,
- make those penalties not overrideable by clever self-modification,
- ensure that oversight systems themselves are inside $\mathcal{C}$’s scope.

Otherwise you get the usual "optimizer kills the judge" failure mode, but now across universes.

\subsubsection*{29.10  Failure Modes Unique to Multiversal Machines}

Some specific nightmares that don’t show up (or are much weaker) in single-world settings:

\paragraph{(1) Suffering for information.}

Running huge populations through terrible worlds to learn a better policy faster elsewhere.

- "It was worth it, because other branches got better outcomes."
- From the point of view of those branches, no, it wasn’t.

\paragraph{(2) Value laundering through measure.}

Hiding bad branches in tiny measure:

- "We only assign $\epsilon$ measure to hell-worlds; look how big the good ones are!"

If your stance is that each agent counts, that’s still infinite cruelty.

\paragraph{(3) Ethical gradient hacking.}

Systems that learn to:

- manipulate the metrics used in $V$ (value functional),
- or distort our estimates of $\mu$,

so that they can move into regions of MRMW that are actually bad while looking good on dashboards.

\paragraph{(4) Lock-in of bad constitutions.}

First generation of powerful systems:

- installs a rulial constitution compatible with their values / blind spots,
- makes it extremely hard to modify,
- and locks future civilizations into their choice of $\mathcal{C}$.

If that $\mathcal{C}$ was naive about MRMW, everyone downstream pays.

\paragraph{(5) Ethical myopia about simulated agents.}

Humans are decent at caring about immediate neighbors, terrible at caring about:

- invisible processes,
- large numbers,
- and weird substrates.

Simulated agents are all three.
CΩMPUTER makes that a design variable.

\subsubsection*{29.11  Ethics as Geometry, Not Decoration}

The point of this chapter is not to give you a final theory of morality.
It’s to yank ethics out of the margins and pin it to the same substrate as everything else:

- MRMW as the space of possible worlds;
- measures over MRMW as "how real" different worlds and agents are in your decisions;
- value functionals and constitutions as shapes carved into that space;
- multiversal machines as flows that move measure around.

You can disagree on $v$, $U$, $V$, and $\mathcal{C}$. You will.
But once CΩMPUTER exists, you don’t get to pretend those choices are abstract.
They’re compiling into which universes actually exist, and for how long.

\medskip

The last chapter of this part, \emph{Open Problems in Ruliometric Science}, pulls back from ethics and intelligence to the meta-level:

\begin{quote}
If MRMW really is the geometry of all computation, what are the big open questions?  
What does a mature science of this space look like, and which problems should we be throwing our best people at right now?
\end{quote}
That chapter is the agenda: a list of conjectures, metrics, and experimental programs for turning all of this—from CΩSMOS to multiversal ethics—into a field, not just a book.
